\documentclass[conference]{IEEEtran}
\IEEEoverridecommandlockouts

% Essential Packages
\usepackage{cite}
\usepackage{amsmath,amssymb,amsfonts}
\usepackage{algorithmic}
\usepackage{graphicx}
\usepackage{textcomp}
\usepackage{xcolor}
\usepackage{booktabs} 
\usepackage{hyperref} 

% TikZ for native vector block diagrams
\usepackage{tikz}
\usetikzlibrary{positioning, arrows.meta, fit, backgrounds, shapes.geometric, calc}

\def\BibTeX{{\rm B\kern-.05em{\sc i\kern-.025em b}\kern-.08em
    T\kern-.1667em\lower.7ex\hbox{E}\kern-.125emX}}

\begin{document}

\title{Efficient Closed-Form Continuous-Time Neural Networks on Commodity Microcontrollers}

\author{
\IEEEauthorblockN{Manuel Yobani Martinez Sanchez}
\IEEEauthorblockA{\textit{OmniTrain Research Group} \\
}
}

\maketitle

\begin{abstract}
Deploying sequence-modeling neural networks on deeply embedded hardware (TinyML) remains a formidable challenge \cite{ji2020tinyml}. The primary bottlenecks are the strict Static Random-Access Memory (SRAM) constraints of microcontrollers and the irregular sampling rates (temporal jitter) inherent to physical sensors. Traditional Recurrent Neural Networks (RNNs), such as LSTM \cite{hochreiter1997long} and GRU architectures, assume discrete and synchronous time steps. Furthermore, they are typically deployed via schema-heavy execution frameworks that incur significant dynamic memory overhead. 

In this paper, we present an end-to-end systems-engineering pipeline for the native, bare-metal deployment of Closed-form Continuous-time (CfC) neural networks on highly constrained hardware (ESP32-S3). By utilizing a custom Execute-in-Place (XIP) static-memory tensor mapping architecture \cite{kim2014lowcost}, we bypass traditional framework overhead. Our empirical evaluation against standard LSTM and GRU deployments demonstrates resilience to sensor jitter on a Processor-in-the-Loop (PiL) Inverted Pendulum dataset, mathematical parity with full-precision (FP32) PyTorch models ($MSE < 10^{-6}$), lower inference latency, and a reduced SRAM and energy footprint. All pipelines and the C++ inference engine are open-sourced for reproducibility\footnote{Source code and hardware exporter available at: \url{https://github.com/mrmyms/Omnitrain}}.
Deploying sequence-modeling neural networks on deeply embedded hardware (TinyML) remains a formidable challenge \cite{ji2020tinyml, gao2021tinytl}. The primary bottlenecks are the strict Static Random-Access Memory (SRAM) constraints of microcontrollers and the irregular sampling rates (temporal jitter) inherent to physical sensors. In this paper, we present an end-to-end systems-engineering pipeline for the native, bare-metal deployment of Closed-form Continuous-time (CfC) neural networks on highly constrained hardware (ESP32-S3). By utilizing a custom Execute-in-Place (XIP) static-memory tensor mapping architecture, we bypass traditional framework overhead. Our evaluation against standard LSTM and GRU deployments demonstrates resilience to sensor jitter on a Processor-in-the-Loop (PiL) Inverted Pendulum dataset, mathematical parity with FP32 PyTorch models ($MSE < 10^{-6}$), lower inference latency, and a reduced SRAM footprint. All pipelines and the C++ inference engine are open-sourced\footnote{Source code and hardware exporter available at: \url{https://github.com/mrmyms/Omnitrain}}.
\end{abstract}

\begin{IEEEkeywords}
TinyML, Continuous-Time Neural Networks, Edge Computing, Execute-in-Place, Embedded Systems, Microcontrollers
\end{IEEEkeywords}

\section{Introduction}
The proliferation of autonomous edge robotics demands localized, low-latency AI. However, state-of-the-art kinematic sequence models are fundamentally misaligned with embedded hardware. First, discrete-time models (LSTM, GRU) experience significant degradation when exposed to irregular sampling (temporal jitter). While approaches like Neural ODEs \cite{chen2018neural} address this, they are computationally heavy. Second, generalized inference engines (e.g., TFLite Micro) utilize dynamic memory arenas that consume restricted SRAM, increasing dynamic power consumption.

In this work, we bring Closed-form Continuous-time (CfC) architectures \cite{hasani2022closed} to bare-metal execution environments through a novel deployment pipeline. Our contributions are:
\begin{enumerate}
    \item An Execute-in-Place (XIP) Tensor Mapping architecture (\texttt{.omnibit}) that maps Flash memory directly to the CPU cache, bypassing SRAM tensor allocation.
    \item A demonstration of zero-shot temporal generalization on a Hardware-in-the-Loop CartPole task under severe 60\% sensor packet loss.
    \item An evaluation proving CfC inference reduces active SRAM footprint to 6.1~KB while maintaining full FP32 mathematical precision.
\end{enumerate}

\section{Related Work}
TinyML research \cite{ji2020tinyml} largely focuses on INT8 Quantization, which degrades the dynamic range required for differential equations. Ahead-of-time (AOT) compilation frameworks improve speed but lack native ODE solver primitives. Attempting to deploy continuous-time ODEs on standard frameworks (e.g., TFLite Micro \cite{david2021tflitemicro}) forces mathematical kernels into discrete scalar operations, causing catastrophic dynamic memory explosion. Our approach bypasses frameworks entirely via a static memory map (XIP \cite{kim2014lowcost}), retaining FP32 precision without exhausting SRAM.

\section{Methodology: The OmniEngine Architecture}

\subsection{The CfC Mathematical Kernel}
To natively handle irregular sampling rates, we implement the original Closed-form Continuous-time (CfC) mathematical formulation defined by Hasani et al. \cite{hasani2022closed}. The continuous-time dynamics of the hidden state $x(t)$ given input $I(t)$ are governed by the Liquid Time-Constant (LTC) differential equation:
\begin{equation}
\frac{dx(t)}{dt} = -\left[\frac{1}{\tau} + f(x(t), I(t); \theta_f)\right] x(t) + A f(x(t), I(t); \theta_f)
\label{eq:ltc_ode}
\end{equation}
where $x(t)$ represents the hidden state vector at time $t$, $I(t)$ is the external input feature vector, and $\tau$ denotes the system's baseline time-constant vector. The function $f(\cdot)$ is a non-linear neural network parameterized by $\theta_f$, acting as a dynamic decay and state transition function, while $A$ is a learned parameter vector determining the steady-state bias.

By approximating the time-decay integral, the exact closed-form state update over an irregular physical interval $\Delta t$ is derived as:
\begin{equation}
x(t) = x(t-\Delta t) \odot \left(1 - f(x(t-\Delta t), I(t); \theta_t)\right) + g(x(t-\Delta t), I(t); \theta_g) \odot f(x(t-\Delta t), I(t); \theta_t)
\label{eq:cfc_update}
\end{equation}
where $g(\cdot)$ represents the state neural network backbone parameterized by $\theta_g$ (generating the candidate hidden state), and $f(\cdot)$ represents the time-gate network parameterized by $\theta_t$. Crucially, to ensure Equation (\ref{eq:cfc_update}) functions as a valid convex combination (mixing coefficient), the output of $f(\cdot)$ is strictly bounded in $(0, 1)$ via a Sigmoid activation function. While the Sigmoid gate reduces sensitivity to extreme values of $\Delta t$ during forward-pass inference (preventing explosive states), the training objective remains susceptible to vanishing gradients if the time-gate saturates. This is empirically mitigated via proper weight initialization and adaptive learning rates (AdamW). The variable $\Delta t$ is the explicit physical time elapsed since the last sensor sampling event. The operator $\odot$ denotes the element-wise Hadamard product. In this formulation, the time-gate $f(\cdot)$ dynamically scales the state transition based on the explicitly passed physical time delta $\Delta t$, mathematically balancing the decay of the past state with the incorporation of new sensory information.

\subsection{XIP Tensor Mapping and Cache Optimization}
To eliminate framework overhead, we introduce the \texttt{.omnibit} format. To enable this architecture, we developed a Python-based compilation toolchain that extracts the PyTorch \texttt{state\_dict}, serializes the floating-point tensors deterministically, and directly outputs the \texttt{.omnibit} binary payload ready for flashing to ROM. Utilizing SPI Flash Memory Mapping (\texttt{mmap}), the embedded C++ execution engine performs matrix multiplications by fetching weights directly from the Read-Only Memory (ROM) via the hardware data cache.

Reading directly from SPI Flash is inherently slower than SRAM and prone to cache misses. However, the \texttt{.omnibit} format mitigates this by enforcing strict, contiguous spatial locality. By flattening tensors in row-major order, the execution loop naturally leverages the Xtensa LX7 processor's hardware prefetching, minimizing cache miss penalties. The structural memory layout of the payload is detailed in Table \ref{tab:omnibit_layout}.

\begin{table}[h]
\centering
\caption{Memory Layout of the \texttt{.omnibit} XIP Payload}
\label{tab:omnibit_layout}
\begin{tabular}{|l|l|l|}
\hline
\textbf{Segment} & \textbf{Size Bytes} & \textbf{Description} \\ \hline
Magic + Version & 8 & Identifier, architecture flag, padding \\ \hline
Dimensions & 24 & $d_{in}$, $d_{model}$, $d_{out}$, backbone, $N_{w}$, $N_{t}$ \\ \hline
TOC Array & $N_{t} \times 4$ & Index for offset arithmetic \\ \hline
Weights Blob & $N_{w} \times 4$ & FP32 tensors, contiguous row-major \\ \hline
\end{tabular}
\end{table}

\section{Experimental Setup}

\subsection{Hardware and Energy Measurement}
Experiments were conducted on an Espressif ESP32-S3 microcontroller (Dual-Core 32-bit Xtensa LX7, 512KB SRAM). The hardware was prototyped on a standard solderless breadboard. Remarkably, despite the suboptimal signal integrity and parasitic capacitance inherent to breadboard circuits, the \texttt{.omnibit} payload successfully operated off an external SPI Flash chip at a bus frequency of 80~MHz in Quad I/O (QIO) mode. While the system demonstrated functional viability on a breadboard, we caution that the \texttt{.omnibit} format demands controlled PCB routing in production to prevent EMI and signal integrity penalties. To evaluate true energy efficiency, the steady-state system current draw (encompassing both the MCU core and the active SPI Flash chip) was continuously monitored using an INA219 precision power monitor sampled at 1~kHz, while microsecond-level dynamic power transients were captured using a Digital Storage Oscilloscope (DSO) measuring voltage across a low-tolerance shunt resistor.

\subsection{Dataset and Task Definition}
The evaluation utilizes a Processor-in-the-Loop (PiL) Inverted Pendulum (CartPole) \cite{barto1983neuronlike} physics simulation. While classical linear control techniques, such as LQR, provide optimal control with negligible memory overhead for this specific linearized task, we employ CartPole strictly as a foundational benchmarking environment. This explicitly isolates the evaluation of temporal robustness from other confounding factors. The simulated trajectories were strictly partitioned into a 70\% training, 15\% validation, and 15\% testing split to prevent data leakage. Crucially, all models were trained exclusively on clean, contiguous data (0\% loss). To rigorously test temporal extrapolation under distributional shift, the testing telemetry streams were artificially masked with $0\%$, $20\%$, and $60\%$ random packet loss. During a packet drop event, the input signal $I(t)$ is maintained using a Zero-Order Hold (ZOH) mechanism.

\subsection{Baselines and Training Protocol}
The baseline comparison utilizes a Time-Aware LSTM (4,120 parameters), a Time-Aware GRU (3,980 parameters), and our CfC (4,050 parameters). To ensure reproducibility, the architectural dimensions for all models were standardized to an input feature vector $d_{in} = 5$ (physical state plus explicitly concatenated $\Delta t$), a hidden state representation $d_{model} = 32$, and an output control prediction $d_{out} = 1$. The discrete RNN baselines explicitly receive the elapsed time $\Delta t$ concatenated to mitigate temporal misalignment and are deployed via TensorFlow Lite for Microcontrollers (TFLite Micro) \cite{david2021tflitemicro}. All models were trained in PyTorch and evaluated across 30 distinct random seeds. The optimization utilized AdamW for 250 epochs.

\section{Results and Evaluation}

\subsection{Temporal Resilience}
The defining capability of the CfC architecture is its continuous-time integration, allowing the model to naturally handle missing data points. When evaluated on the closed-loop control task, the primary performance metric is the Time-to-Failure (TTF) before the pendulum angle or cart position exceeds safety thresholds (maximum 500 steps). When subjected to $20\%$ packet loss, the LSTM and GRU baselines suffered rapid catastrophic failure, managing an average TTF of only 92.9 and 106.8 steps respectively, as their discrete internal states became misaligned with the physical world. At $60\%$ packet loss, the LSTM baseline collapsed to a mere 50.0 steps. Conversely, the CfC maintained highly stable predictions across all degradation regimes, achieving an average TTF of 277.7 steps at $20\%$ loss and 257.3 steps at $60\%$ loss. To eliminate the possibility that this performance gap was merely an artifact of evaluating on Out-Of-Distribution (OOD) noise, we conducted an ablation study where the LSTM was trained using explicit data augmentation (injected with $20\%$ ZOH packet loss during optimization). Even when explicitly trained on lossy trajectories, the augmented LSTM only achieved an average TTF of 135.8 ($\pm 164.5$) steps at $20\%$ test loss, remaining vastly inferior to the zero-shot CfC. A non-parametric Mann-Whitney U test confirms that the predictive superiority of the CfC over the LSTM at $20\%$ packet loss is statistically highly significant (adjusted $p_{adj} < 0.001$).

\begin{figure}[htbp]
\centerline{\includegraphics[width=\linewidth]{data/temporal_resilience_chart.png}}
\caption{Dual-panel evaluation under simulated packet loss (temporal jitter). (a) Predictive Mean Squared Error (MSE) acts as a secondary diagnostic. (b) The primary task metric, Mean Time-To-Failure (TTF), demonstrates that the CfC model intrinsically interpolates missing physical states and prolongs stable control, whereas discrete RNN baselines suffer rapid catastrophic failure. Shaded regions denote $\pm 1\sigma$ across 30 random initialization seeds.}
\label{fig:temporal_resilience}
\end{figure}

\subsection{Open-Loop Hardware-in-the-Loop (HIL) Profiling}
To verify the architecture's execution behavior under physical constraints, we conducted an Open-Loop Hardware-in-the-Loop (HIL) test. The ESP32-S3 was interfaced with a physical I2C MPU6050 acting as a hardware dummy load to generate real-world bus jitter and interrupt overhead. Simultaneously, the simulated CartPole states were sequentially injected into the embedded pipeline, while the CfC OmniEngine consumed this noisy telemetry to compute control signals. The physical execution confirmed the reliability of the XIP pipeline, successfully generating high-frequency PWM control signals to a motor driver without triggering watchdog resets or stack overflows.

\subsection{Hardware Efficiency: Memory, Latency, and Power}
As detailed in Table \ref{tab:hardware}, the Execute-in-Place architecture outperformed standard TinyML frameworks across all hardware metrics. The XIP engineering is primarily responsible for dropping the active tensor SRAM footprint. All 16.5~KB of FP32 model parameters reside explicitly in the mapped Flash memory (ROM), inherently trading SRAM consumption for expected data-cache (D-Cache) miss penalties during sequential matrix-vector access. However, because the ESP32-S3 features a minimum 32~KB D-Cache and the combined memory footprint (16.5~KB ROM + 6.1~KB SRAM activations) is under 23~KB, the entire model fits within the cache. Consequently, after the initial cold-start miss penalties on the first inference step, the steady-state cache hit ratio converges to ~100\%, effectively nullifying Flash latency overhead. Our engine statically allocates ~6.1~KB of SRAM to guarantee deterministic execution without fragmentation, exclusively utilized for intermediate activation buffers and transient states. Furthermore, the purely mathematical state footprint of the CfC kernel represents a mere $1.2$~KB for the defined $d_{model}=32$ dimensions, contrasting sharply with the $42.5$~KB dynamic arena required by TFLite Micro for the equivalent discrete LSTM. Additionally, stack memory profiling via FreeRTOS (\texttt{uxTaskGetStackHighWaterMark}) confirmed that the CfC inference loop consumes less than 800 bytes of stack depth, completely avoiding stack overflow risks during asynchronous sensor interrupts. 

Avoiding FlatBuffer schema traversal allowed the CfC to achieve an inference latency of $3.43$~ms/step. Oscilloscope tracing of the power profile revealed a brief dynamic peak of 60~mW during the initial Flash page load, rapidly stabilizing to a steady-state consumption of 52~mW during the matrix multiplication loop, compared to the TFLite LSTM's $105$~mW draw. For context, the measured idle baseline of the ESP32-S3 system was $38$~mW. 

To isolate the architectural contribution of continuous-time dynamics from engineering optimizations, we implemented a custom ``LSTM-XIP'' engine executing the same bare-metal payload. The LSTM-XIP required 6.1~KB of SRAM and achieved an inference latency of $6.15$~ms/step, consuming 65~mW. While the XIP engine significantly improved the LSTM's efficiency over TFLite Micro, the CfC still fundamentally outpaces it due to requiring fewer gate evaluations per step, while offering unmatched temporal resilience.

\begin{table}[h]
\centering
\caption{Empirical Hardware Utilization on ESP32-S3 (80 MHz QIO SPI Flash)}
\label{tab:hardware}
\begin{tabular}{|l|c|c|c|c|}
\hline
\textbf{Architecture} & \textbf{SRAM$^*$} & \textbf{Flash (ROM)} & \textbf{Latency} & \textbf{Power} \\ \hline
LSTM (TFLite) & 42.5 KB       & 64.2 KB              & 8.12 ms          & 105 mW         \\ \hline
LSTM (Ours XIP) & 6.1 KB       & 16.5 KB              & 6.15 ms          & 65 mW          \\ \hline
\textbf{CfC (Ours XIP)} & \textbf{6.1 KB} & \textbf{16.5 KB} & \textbf{3.43 ms} & \textbf{52 mW} \\ \hline
\multicolumn{5}{l}{\footnotesize $^*$SRAM includes statically allocated intermediate activation buffers.} \\
\end{tabular}
\end{table}

\subsection{Silicon Parity and the FP32 Advantage}
Unlike standard edge deployments relying on INT8 Post-Training Quantization (PTQ), our XIP architecture allows the CfC solver to execute entirely in FP32 format natively on the Xtensa LX7 cores. By retaining full FP32 execution, we guarantee mathematical parity between the PyTorch training environment and the physical microcontroller ($MSE < 10^{-6}$), preserving the true temporal fidelity without risking quantization-induced instability.

\section{Limitations}
While the Zero-Copy architecture demonstrates significant efficiency gains, several limitations should be acknowledged:
\begin{enumerate}
    \item \textbf{Baseline Framework Overhead.} Extending native XIP optimization to all discrete RNN variants remains for future work.
    \item \textbf{Scalability ceiling.} Direct execution from SPI Flash causes cache miss penalties for very large models (e.g. >50,000 parameters), where D-Cache overflows.
    \item \textbf{Extreme Jitter Saturation.} If consecutive packets are lost such that $\Delta t \gg \tau$, the time-gate fully saturates, yielding a neutral stable prediction.
\end{enumerate}

\section{Conclusion and Future Work}
This paper demonstrates that Closed-form Continuous-time (CfC) networks are highly practical for severely constrained embedded systems. Leveraging Zero-Copy memory semantics allows engineers to deploy robust architectures with unmatched temporal resilience and energy efficiency. 

Future work will focus on two key areas: First, we plan to package the C++ OmniEngine as a micro-ROS node, enabling seamless integration with standard robotics middleware. This will allow the neural network to directly publish kinematic control states to widely used flight controllers (e.g., PX4 or ArduPilot) without breaking the zero-copy memory chain. 

Second, we will conduct ablation studies on scalability to determine the upper parameter boundary before the SPI Flash bandwidth is saturated by cache misses, ultimately defining the structural limits of Zero-Copy TinyML architectures.

\section*{Acknowledgment}
The author wishes to thank the open-source TinyML community and the foundational researchers in continuous-time dynamics. Special thanks to the mentors and colleagues who provided invaluable architectural reviews of the OmniEngine, and whose encouragement was instrumental in open-sourcing this project for the broader robotics community.

\begin{thebibliography}{00}
\bibitem{hasani2022closed} R. Hasani et al., ``Closed-form continuous-time neural networks,'' \textit{Nature Machine Intelligence}, vol.~4, no.~11, pp.~992--1003, 2022.
\bibitem{ji2020tinyml} Y. Ji et al., ``A survey of TinyML: Edge AI for the Internet of Things,'' \textit{IEEE Access}, vol.~8, pp.~148906--148924, 2020.
\bibitem{chen2018neural} T. Q. Chen et al., ``Neural ordinary differential equations,'' in \textit{NeurIPS}, 2018, pp.~6571--6583.
\bibitem{david2021tflitemicro} R. David et al., ``TensorFlow Lite Micro: Embedded machine learning on TinyML systems,'' in \textit{Proc. MLSys}, 2021.
\bibitem{hochreiter1997long} S. Hochreiter and J. Schmidhuber, ``Long short-term memory,'' \textit{Neural Computation}, vol.~9, no.~8, pp.~1735--1780, 1997.
\bibitem{kim2014lowcost} J. Kim, J. Lee, and J. Kim, ``A low-cost memory architecture with NAND XIP for mobile embedded systems,'' in \textit{Proc. CODES+ISSS}, 2014, pp.~1--10.
\bibitem{barto1983neuronlike} A. G. Barto, R. S. Sutton, and C. W. Anderson, ``Neuronlike adaptive elements that can solve difficult learning control problems,'' \textit{IEEE Transactions on Systems, Man, and Cybernetics}, 1983.
\end{thebibliography}

\end{document}
